%==============================================================================
\section{Public Identity and Access}
\label{sec:public-identity}
%==============================================================================
Some people wish their identity to be publicly known, the way to access them publicly known, or both. To support this, we define two publishable self-certifying records---identity and access---that a person can post or disseminate outside the grassroots platform, e.g. on a personal website, in an email signature, through social media, or via a dissemination protocol such as Nostr or Scuttlebutt.
The two records are independent: a person may publish either, both, or neither.  Each record can be realized as an executable URL: when opened, the grassroots app adds the person's identity or access information directly.

 
%------------------------------------------------------------------------------
\subsection{Identity Record}
\label{sec:identity-record}
%------------------------------------------------------------------------------
An \emph{identity record} binds a person's public key to public information about them.  It consists of:
\begin{itemize}
\item the public key of the person (natural or legal);
\item the public keys of the person's identity custodians and the supermajority threshold required for identity replacement;
\item any information the person wishes to make public (e.g., first and last name, organizational affiliation);
\item if this record replaces or updates a prior identity record: a hash of the public record of the duly-signed identity replacement protocol by the identity custodians of the previous identity;
\item a signature by the private key of the person.
\end{itemize}

Identity custodians, the supermajority threshold, and the identity replacement protocol are defined in a companion paper.%TODO: add citation to the paper that defines identity custodians and the replacement protocol.

%------------------------------------------------------------------------------
\subsection{Access Record}
\label{sec:access-record}
%------------------------------------------------------------------------------
An \emph{access record} binds a person's public key to the rendezvous servers they use.  It consists of:
\begin{itemize}
\item the public key of the person;
\item the addresses of the rendezvous servers the person uses;
\item an epoch number;
\item a signature by the private key of the person.
\end{itemize}

Access records carry no expiration.  A later record supersedes an earlier one by epoch number: a recipient holding multiple records signed by the same public key retains only the record with the highest epoch number.  To change the list of rendezvous servers, the person publishes a new access record with a higher epoch number.

%------------------------------------------------------------------------------
\subsection{Publication Choices}
\label{sec:publication-choices}
%------------------------------------------------------------------------------
Publishing only an identity record makes a person publicly identifiable but exposes no information about where they can be reached.  Publishing only an access record makes a person reachable via the rendezvous mechanism of Section~\ref{sec:rendezvous} to anyone who knows their public key, without revealing any additional information about them.  Publishing both makes the person both publicly identifiable and reachable.  Publishing neither keeps the person publicly unidentifiable and publicly inaccessible; they remain reachable through friend-mediated introduction, BLE adjacency, and private sharing of their public key and rendezvous servers.
